% source: https://github.com/pfradin/luadraw/discussions/249#discussioncomment-16927920 (pfradin, block 1)
\documentclass{standalone}%compile with lualatex only
\usepackage[svgnames]{xcolor}
\usepackage{luadraw}%https://github.com/pfradin/luadraw/
\usepackage{fourier-otf}

\begin{luacode*}
local ld= luadraw
local cpx = ld.cpx
local Z = cpx.Z

function ld.parabola(A,B,C)
    A = cpx.toComplex(A); B = cpx.toComplex(B); C = cpx.toComplex(C)
    local x1, x2, x3 = A.re, B.re, C.re -- the x-coordinates must be different
    local y1,y2,y3 = A.im, B.im, C.im
    local f = function(t)
        return (t-x2)*(t-x3)/(x1-x2)/(x1-x3)*y1 + (t-x1)*(t-x3)/(x2-x1)/(x2-x3)*y2 + (t-x2)*(t-x1)/(x3-x2)/(x3-x1)*y3
    end
    return f
end
\end{luacode*}

\begin{document}

\begin{luadraw}{name=parabola}
local ld = luadraw
local cpx = ld.cpx
local Z = cpx.Z

local g = ld.graph:new{size={10,10}}
local A, B, C = Z(-3,2), Z(0,1), Z(2,-2)
g:Daxes({0,1,1}, {grid=true, arrows="-Stealth", legend={"$x$","$y$"}})
local f = ld.parabola(A,B,C)
g:Ddots({A,B,C})
g:Dcartesian(f, {draw_options="red, line width=0.8pt"}); 
g:Show()
\end{luadraw}
\end{document} 
